\documentclass[11pt]{article}

\usepackage[a4paper, margin=2.5cm]{geometry}
\usepackage{graphicx}
\usepackage{booktabs}
\usepackage{hyperref}
\usepackage{amsmath}

\title{Unsupervised Learning Project:\\
A Reproducible Experimental Study of Clustering}
\author{
Alexandre Santos (72970) \\
David Natal (72997) \\
Miguel Mestre (73018)
}
\date{30 April 2026}

\begin{document}

\maketitle

\section*{Front Matter}
\textbf{Project title:} Hotel Booking Demand Clustering\\
\textbf{Repository:} [Insert GitHub URL]\\
\textbf{Version:} [Insert commit hash]

\section{RQ1 and Segmentation Time}

\textbf{Research Question:}\\
Can hotel booking records be partitioned into distinct, interpretable behavioural segments
using only pre-stay observable characteristics --- daily rate, total estimated spend,
length of stay, booking horizon, meal plan, deposit commitment, special requests, and
parking need --- to identify guest profiles that support targeted revenue management and
personalised service strategies?

\textbf{Unit of analysis:}\\
One row corresponds to one booking record. A single guest may appear multiple times
across different stays.

\textbf{Segmentation time:}\\
Features are restricted to information available at booking time. Post-event fields
(\texttt{is\_canceled}, \texttt{reservation\_status}, \texttt{reservation\_status\_date})
are excluded from clustering and retained only for post-hoc profiling.

\subsection*{Cluster Findings ($k = 3$)}

Running K-Means with $k = 3$ on the pre-stay feature set yields three stable and
interpretable segments (ARI $= 0.9998$ across 10 random seeds):

\begin{center}
\begin{tabular}{clrp{7cm}}
\toprule
Cluster & Label & Size & Key characteristics \\
\midrule
1 & \textbf{Low-Cost Guest} & 48{,}475 &
    Lowest total spend (log spend $-45$\% vs.\ mean), shortest stays (1.9 nights),
    last-minute booking ($-73$\% lead-time rank), highest car-parking need (+48\%).
    Budget, spontaneous city traveller. \\
0 & \textbf{Medium-Value Guest} & 53{,}403 &
    Near-average ADR (\texteuro102), moderate stay length (3.3 nights),
    non-refundable deposit (+73\% vs.\ mean), planned ahead (+55\% lead-time rank).
    Standard urban booker with commitment. \\
2 & \textbf{Premium Guest} & 17{,}511 &
    Highest total spend (+136\% vs.\ mean), longest stays (8.1 nights),
    resort hotel dominant, above-average special requests (+23\%).
    High-value leisure segment. \\
\bottomrule
\end{tabular}
\end{center}

These three profiles directly address the research question: Low-Cost, Medium-Value, and
Premium guests exhibit clearly distinct spending patterns, stay durations, and booking
behaviours --- segments with direct implications for dynamic pricing, upselling, and
service allocation.

\section{Dataset Version Control and Documentation}

\textbf{Dataset version:}\\
\texttt{hotel\_bookings\_course\_release\_v1.csv}

\textbf{Fingerprint (SHA-256):}\\
\texttt{7c2ae42a7353905ea136e5c2287f17c92c5435826598bfbb8491c6f0c7b1fc06}

\textbf{Source and Usage Terms:}\\
Hotel PMS data from Portugal. [Insert usage terms or reference note here]

\textbf{Known issues:}
\begin{itemize}
    \item High missingness: \texttt{company} ($\approx$94\%), \texttt{agent} ($\approx$13.5\%)
    \item Minor missingness: \texttt{country}, \texttt{children}
    \item Strong skew and outliers: \texttt{adr}, \texttt{lead\_time}
    \item Zero-ADR rows ($\approx$1{,}960): complementary or erroneous entries, excluded
          from clustering (see Section~\ref{sec:representation})
\end{itemize}

\section{Exploratory Data Analysis}

\textbf{Missingness:}
\begin{itemize}
    \item Numerical: only \texttt{children} has missing values (4 rows)
    \item Categorical: \texttt{company} (112{,}593), \texttt{agent} (16{,}340),
          \texttt{country} (488)
\end{itemize}

\textbf{Outliers:}
\begin{itemize}
    \item \texttt{adr}: heavy right tail (values up to \texteuro5{,}400); one negative
          entry clipped to zero and subsequently excluded
    \item \texttt{lead\_time}: strong variability (0--737 days); binned into quantile
          ranks to reduce influence
\end{itemize}

\textbf{Dropped variables (pre-processing):}
\begin{itemize}
    \item \texttt{company}, \texttt{agent}: high-cardinality ID-like fields with no
          principled encoding strategy
\end{itemize}

\textbf{Leakage exclusions (clustering only):}
\begin{itemize}
    \item \texttt{is\_canceled}: post-decision outcome
    \item \texttt{reservation\_status}, \texttt{reservation\_status\_date}: post-arrival
          fields
\end{itemize}

\section{Main Representation (RQ2)}\label{sec:representation}

\textbf{Clustering dataset:} 117{,}430 bookings after removing zero- and negative-ADR
rows (\texttt{adr} $> 0$). Zero-ADR records correspond to complimentary stays and
artefact entries that otherwise form a spurious singleton cluster.

\subsection*{Included features}

\begin{center}
\begin{tabular}{lll}
\toprule
Feature & Type & Scaling / Encoding \\
\midrule
\texttt{total\_spend\_log} & Continuous & RobustScaler \\
\texttt{total\_nights}     & Continuous & RobustScaler \\
\texttt{is\_non\_refund}   & Binary     & Raw $\times 3.0$ \\
\texttt{total\_of\_special\_requests} & Count & Raw \\
\texttt{required\_car\_parking\_spaces} & Binary & Raw \\
\texttt{meal\_encoded}     & Ordinal (0--3) & Raw $\times 0.5$ \\
\texttt{lead\_time\_rank}  & Ordinal (0--2) & Raw \\
\bottomrule
\end{tabular}
\end{center}

\subsection*{Derived features and preprocessing decisions}

\begin{itemize}
    \item \textbf{\texttt{total\_spend\_log}}: computed as
          $\log(1 + \texttt{adr}) \times \texttt{total\_nights}$.
          Multiplying the log-transformed rate by duration produces a compressed spend
          proxy that rewards both price level and length of stay.
          \texttt{adr\_log} is \emph{not} included as a separate feature because it is
          collinear with \texttt{total\_spend\_log} (spend $\approx$ rate $\times$
          nights), which would create an elongated manifold K-Means handles poorly.

    \item \textbf{\texttt{total\_nights}}: sum of weeknight and weekend-night stays.

    \item \textbf{\texttt{is\_non\_refund}}: binary indicator
          (\texttt{deposit\_type} $=$ ``Non Refund'' $\rightarrow 1$, otherwise $0$).
          Replaces a three-level ordinal encoding that artificially inflated Euclidean
          distances. Upweighted by $\times 3.0$ to ensure deposit commitment is a
          primary axis of differentiation.

    \item \textbf{\texttt{lead\_time\_rank}}: \texttt{lead\_time} binned into three
          equal-population quantile groups via \texttt{pd.qcut} with \texttt{q=3}
          (0 = last-minute: 0--32 days; 1 = mid: 33--124 days;
           2 = early planner: 125--737 days).
          Binning reduces right-tail dominance while preserving the booking-horizon
          ordering.

    \item \textbf{\texttt{meal\_encoded}}: ordinal meal-plan encoding
          (SC/Undefined $= 0$, BB $= 1$, HB $= 2$, FB $= 3$), downweighted by
          $\times 0.5$ to prevent a coarse categorical proxy from dominating distances.

    \item \textbf{RobustScaler} (IQR-based) applied to \texttt{total\_spend\_log} and
          \texttt{total\_nights} to limit the influence of extreme outliers.
          Binary and ordinal features are left on their natural scales before weighting.
\end{itemize}

\textbf{Excluded variables:}
\begin{itemize}
    \item \texttt{country}: retained for post-hoc profiling only; too high a cardinality
          for direct inclusion
    \item \texttt{adr\_log}: intermediate computation only; redundant with
          \texttt{total\_spend\_log}
\end{itemize}

\textbf{Representation ID:} \texttt{v1\_robust\_weighted}

\section{Baseline Modelling Evidence}

\textbf{Fixed $K$ values considered:} $K = 2, 3, 4$ \\
\textbf{Justification:} The elbow plot (inertia vs.\ $K$) flattens after $K = 3$.
Silhouette peaks at $K = 2$ and $K = 4$; $K = 3$ is selected as the primary solution
because it balances compactness with business interpretability (three actionable segments)
and achieves the highest ARI stability (0.9998).

\subsection*{K-Means Results \normalfont\small(\textbf{Runtime:} 2 min 27 s on a 15{,}000-row sample)}

All runs: \texttt{n\_init=10}, \texttt{random\_state=42}, feature matrix
\texttt{v1\_robust\_weighted}.

\begin{center}
\begin{tabular}{lccc}
\toprule
$K$ & Silhouette $\uparrow$ & Calinski-Harabasz $\uparrow$ & Davies-Bouldin $\downarrow$ \\
\midrule
2 & 0.454 & 74{,}227 & 0.969 \\
3 & 0.276 & 59{,}877 & 1.296 \\
4 & 0.294 & 59{,}549 & 1.151 \\
5 & 0.312 & 59{,}371 & 1.044 \\
\bottomrule
\end{tabular}
\end{center}

\textbf{Observation:}
$K = 2$ maximises the silhouette but merges the two short-stay segments into one, losing
interpretability. $K = 3$ produces three semantically distinct profiles with excellent
ARI stability ($0.9998 \pm 0.0002$). $K = 4$ splits the high-value resort segment
further but with moderate ARI stability ($0.995$).

\subsection*{iK-Means Results \normalfont\small(\textbf{Runtime:} 8 min 30 s on a 15{,}000-row sample)}

iK-Means (Mirkin 2005) was run on the full 117{,}430-row matrix to obtain an
initialisation-free $K$ estimate.

\begin{itemize}
    \item Anomalous clusters extracted: 12
    \item Final $k$ (clusters $\geq 1{,}000$ members): $k = 8$
    \item Silhouette: 0.311
    \item Calinski-Harabasz: 53{,}163
    \item Davies-Bouldin: 0.995
\end{itemize}

\textbf{Observation:}
iK-Means confirms that more than 3 structure-bearing subgroups exist in the data, but the
additional clusters fragment the three main segments identified by K-Means rather than
revealing qualitatively new profiles.

\subsection*{Agglomerative Hierarchical Clustering (Ward Linkage)}

Run on a reproducible subsample of $n = 30{,}000$ (seed 42).

\begin{center}
\begin{tabular}{lccc}
\toprule
$K$ & Silhouette $\uparrow$ & Calinski-Harabasz $\uparrow$ & Davies-Bouldin $\downarrow$ \\
\midrule
2 & 0.478 & 16{,}763 & 0.911 \\
3 & 0.203 & 12{,}354 & 1.554 \\
4 & 0.234 & 11{,}194 & 1.287 \\
\bottomrule
\end{tabular}
\end{center}

\textbf{Observation:}
Hierarchical clustering corroborates the two- and four-cluster structure found by K-Means.
The Ward dendrogram shows a clear first split at $k = 2$ and a secondary split producing
a high-spend long-stay group, consistent with Cluster 2 from K-Means.

\subsection*{Cluster Profile Summary ($k = 3$, K-Means)}

\begin{center}
\begin{tabular}{lrrrr}
\toprule
 & Cluster 0 & Cluster 1 & Cluster 2 & Grand mean \\
\midrule
Size                          & 53{,}403 & 48{,}475 & 17{,}511 & --- \\
ADR (\texteuro)               & 102.1  & 98.7   & 109.8  & 102.7 \\
Total nights                  & 3.3    & 1.9    & 8.1    & 3.4   \\
Total spend (log)             & 14.9   & 8.6    & 36.7   & 15.5  \\
\texttt{is\_non\_refund}      & 0.21   & 0.06   & 0.04   & 0.12  \\
Lead time rank                & 1.55   & 0.27   & 1.34   & 1.07  \\
Special requests              & 0.52   & 0.58   & 0.70   & 0.60  \\
Car parking                   & 0.03   & 0.09   & 0.06   & 0.06  \\
\bottomrule
\end{tabular}
\end{center}

\section{Sensitivity and Robustness}

\subsection*{Stability Analysis (ARI across 10 seeds)}

\begin{center}
\begin{tabular}{lcc}
\toprule
$K$ & Mean ARI $\uparrow$ & Std ARI \\
\midrule
2 & 0.9999 & 0.0001 \\
3 & 0.9998 & 0.0002 \\
4 & 0.9948 & 0.0045 \\
5 & 0.9707 & 0.0580 \\
\bottomrule
\end{tabular}
\end{center}

$K = 3$ achieves near-perfect cross-seed agreement, confirming that the three-cluster
solution is not an artefact of random initialisation.

\subsection*{Scaler Sensitivity (uniform scalers, $k = 4$)}

\begin{center}
\begin{tabular}{lccc}
\toprule
Scaler & Silhouette $\uparrow$ & Calinski-Harabasz $\uparrow$ & Davies-Bouldin $\downarrow$ \\
\midrule
Standard (uniform) & 0.398 & 43{,}698 & 0.980 \\
Robust (uniform)   & 0.325 & 68{,}940 & 1.039 \\
MinMax (uniform)   & \textbf{0.609} & \textbf{246{,}815} & \textbf{0.622} \\
\bottomrule
\end{tabular}
\end{center}

\textbf{Observation:}
MinMax scaling inflates silhouette and Calinski-Harabasz scores substantially because it
compresses all features to $[0, 1]$, producing artificially tight clusters dominated by
binary features. The mixed RobustScaler pipeline used in the main analysis is more
appropriate for this dataset's skewed continuous features.

\section{Experiment Logging and Repository Status}

\textbf{Logging file:} \texttt{experiments.csv}

\textbf{Schema:}
\begin{center}
\begin{tabular}{ll}
\toprule
Column & Description \\
\midrule
date/run\_id       & Execution timestamp \\
representation\_id & Feature set identifier \\
model             & Algorithm used \\
parameters        & Model hyperparameters \\
seed              & Random seed \\
silhouette        & Silhouette score \\
calinski          & Calinski-Harabasz index \\
davies            & Davies-Bouldin index \\
diagnostics       & OK / error flag \\
notes             & Free-text observations \\
\bottomrule
\end{tabular}
\end{center}

\textbf{Environment:} Controlled via \texttt{environment.yml}\\
\textbf{Pipeline:} Fully reproducible; re-run from cell \texttt{2cc7ddee} in
\texttt{project.ipynb} to regenerate all results.

\newpage
\section*{Appendix (Suggested)}
\begin{itemize}
    \item Feature-selection table
    \item RobustScaler / weighting pipeline sketch
    \item Missingness table and outlier summary
    \item Elbow plot (inertia vs.\ $K$)
    \item Validity-index comparison plot (Silhouette, CH, DB across algorithms)
    \item Cluster $\Delta\%$ deviation table ($k = 3$)
    \item Intra-cluster variance table ($k = 2, 3, 4$)
    \item \texttt{experiments.csv} excerpt
\end{itemize}

\end{document}
